\chapter{李群与李代数}
\label{ch:lie}

% ════════════════════════════════════════════════════════════════
%  黄金范本章：吸收 视觉SLAM十四讲 ch4 + Barfoot SER2 ch7/8 + SLAM Handbook ch2
%  + Solà micro-Lie(1812.01537)，自包含独立记录，右扰动为主。
%  教学层遵循 v5.0 理论教学规范（动机→反面→历史→理论→陷阱→练习 + 符号表/速查/故障排查）。
% ════════════════════════════════════════════════════════════════

\begin{practice}[前置自测]
开始之前，先试答下列问题。若已能流畅作答，可只速读本章表格与速查卡；若卡住，正是本章要为你解决的。
\begin{enumerate}
  \item 两个旋转矩阵 $\mathbf{R}_1,\mathbf{R}_2\in\SO(3)$ 相加，结果还是旋转矩阵吗？相乘呢？
  \item 优化一个相机位姿 $\mathbf{T}$ 时，为什么不能像优化普通向量那样对 $\mathbf{T}$ 的 $16$ 个元素直接做梯度下降？
  \item 旋转向量（轴角）$\boldsymbol{\phi}=\theta\mathbf{a}$ 与旋转矩阵 $\mathbf{R}$ 之间的换算公式叫什么？它和“矩阵指数”有何关系？
  \item 给定误差函数 $\mathbf{e}(\mathbf{R})=\mathbf{R}\mathbf{p}-\mathbf{z}$，你打算如何定义并计算 $\partial\mathbf{e}/\partial\mathbf{R}$？“对一个矩阵求导”在这里合法吗？
  \item 在 SLAM 代码里把位姿更新写成 $\mathbf{R}\leftarrow\mathbf{R}\,\mathrm{Exp}(\delta\boldsymbol{\phi})$ 与 $\mathbf{R}\leftarrow\mathrm{Exp}(\delta\boldsymbol{\phi})\,\mathbf{R}$，二者有何区别？哪个是“在机体系下扰动”？
\end{enumerate}
\end{practice}

\section*{本章目标}
学完本章，你应当能够：
\begin{enumerate}
  \item \textbf{说清}旋转/位姿为何构成\emph{群}而非向量空间，以及这给优化带来的根本困难；
  \item \textbf{推导} $\SO(3)$、$\SE(3)$ 的李代数 $\so(3)$、$\se(3)$，并解释“李代数 = 单位元处切空间”的含义；
  \item \textbf{掌握}指数/对数映射（罗德里格斯公式、$\SE(3)$ 指数与左雅可比 $\mathbf{J}$），并理解 $\boldsymbol{\rho}\neq\mathbf{t}$ 的来由；
  \item \textbf{运用} BCH 线性近似与左/右雅可比 $\mathbf{J}_l,\mathbf{J}_r$，在李代数上做加减与微分；
  \item \textbf{独立完成}右扰动模型下旋转点、位姿作用点对扰动的求导，并说明它为何优于直接李代数求导；
  \item \textbf{使用}伴随 $\mathrm{Ad}(\mathbf{T})$ 在左右扰动、不同坐标系之间精确转换；
  \item \textbf{联系}流形 / retraction / $\boxplus\,\boxminus$ 的现代语言，理解右扰动如何对接因子图与 GTSAM/Ceres 求解器；
  \item \textbf{初步处理}李群上的高斯不确定度与位姿复合的协方差传播。
\end{enumerate}

\subsection*{本章知识导航}
本章是一条“\emph{把带约束的旋转/位姿，化为可自由加减、可微分的向量}”的主线。结构如下：

\begin{center}
\small
\begin{tikzpicture}[
  node distance=4mm and 6mm, >=Stealth,
  blk/.style={draw, rounded corners, align=center, font=\small, inner sep=3pt, fill=main!6, draw=main!60!black},
  grp/.style={draw, rounded corners, align=center, font=\small, inner sep=3pt, fill=second!6, draw=second!60!black},
  ln/.style={->, thick, gray!60!black}]
\node[blk] (group) {群 / 矩阵李群\\ $\SO(3),\SE(3)$};
\node[blk, right=of group] (alg) {李代数\\ $\so(3),\se(3)$};
\node[blk, right=of alg] (exp) {指数 / 对数\\ Rodrigues, $\mathbf{J}$};
\node[grp, below=8mm of group] (bch) {BCH\\ $\mathbf{J}_l,\mathbf{J}_r$};
\node[grp, right=of bch] (pert) {扰动模型\\（右扰动为主）};
\node[grp, right=of pert] (adj) {伴随\\ $\mathrm{Ad},\mathrm{ad}$};
\node[blk, below=8mm of bch, fill=third!8, draw=third!60!black] (man) {流形 / retraction\\ $\boxplus\,\boxminus$（对接因子图）};
\node[blk, right=of man, fill=third!8, draw=third!60!black] (unc) {李群上的\\不确定度};
\draw[ln] (group)--(alg); \draw[ln] (alg)--(exp);
\draw[ln] (exp.south) -- ++(0,-4.5mm) -| (bch.north);
\draw[ln] (bch)--(pert); \draw[ln] (pert)--(adj);
\draw[ln] (pert.south)-- (man.north);
\draw[ln] (man)--(unc);
\end{tikzpicture}
\end{center}

\noindent\textbf{推荐路径}：第 1--4 节是任何读者的主干（建立“群—代数—指数”三角）；第 5--7 节是\emph{优化所需}的核心工具（BCH、扰动、伴随），SLAM 后端、VIO、标定章节会反复引用；第 8 节用现代流形语言把前面统一，并说明本书\emph{右扰动}约定的工程归宿（因子图）；第 9 节（不确定度）面向状态估计，可在读完滤波章后回看。带 $*$ 的 Sim(3) 与不确定度四阶项为进阶/选读。

\subsection*{前置知识桥接}
\cref{ch:rigid_body} 已给出旋转的多种\emph{表示}——旋转矩阵 $\mathbf{R}$、旋转向量（轴角）$\boldsymbol{\phi}=\theta\mathbf{a}$、欧拉角、四元数，以及罗德里格斯公式 $\mathbf{R}=\cos\theta\,\mathbf{I}+(1-\cos\theta)\mathbf{a}\mathbf{a}^\top+\sin\theta\,\mathbf{a}^\wedge$ 的“正向”用法。本章要解决的是上一章\emph{回避}的问题：当 $\mathbf{R},\mathbf{T}$ 是\textbf{待估计、待优化}的未知量时，如何对它们求导、做迭代更新。换言之，从“如何表示旋转”进到“\emph{如何估计旋转}”。

\subsection*{如果跳过本章会怎样}
\begin{itemize}
  \item 在 \cref{ch:nlopt}（非线性优化）里，你会写出形如 $\min_{\mathbf{T}}\sum_i\|\mathbf{z}_i-\mathbf{T}\mathbf{p}_i\|^2$ 的问题，却\emph{无法}计算 $\partial/\partial\mathbf{T}$——因为 $\SE(3)$ 没有加法，求导定义直接失效；
  \item 在 VIO / 位姿图优化里，你会发现 Ceres、GTSAM 都要求你提供一个“流形局部参数化（local parameterization / manifold）”，而不解释清楚它就是本章的\emph{右扰动 + 指数映射}，你将只能照抄而不会调试。
\end{itemize}

\begin{table}[H]\centering\small
\caption{本章阅读方式建议}
\begin{tabular}{lll}
\toprule
阅读方式 & 时间 & 适合谁 \\
\midrule
精读（含全部推导与练习） & 5--7 小时 & 首次系统学习、要做后端/标定推导 \\
速读（跳过冗长推导细节） & 2 小时 & 有李群基础、想对齐本书记号与右扰动约定 \\
速查（只看小结的符号表 + 定理速查表） & 15 分钟 & 写代码时回来查公式 \\
\bottomrule
\end{tabular}
\end{table}

\begin{note}
\textbf{本章记号与约定.}\;
旋转矩阵记 $\mathbf{R}\in\SO(3)$（Barfoot\cite{barfoot2024state} 用 $\mathbf{C}$，本书统一 $\mathbf{R}$，与十四讲\cite{gaoxiang2019slam14}、Eigen 一致）；位姿 $\mathbf{T}\in\SE(3)$。
$\so(3)$ 元素 $\boldsymbol{\phi}\in\mathbb{R}^3$，经反对称算子 $(\cdot)^\wedge$ 化为 $\boldsymbol{\phi}^\wedge\in\mathbb{R}^{3\times3}$，逆为 $(\cdot)^\vee$。
$\se(3)$ 元素 $\boldsymbol{\xi}=[\boldsymbol{\rho}^\top,\boldsymbol{\phi}^\top]^\top\in\mathbb{R}^6$（\textbf{平移分量 $\boldsymbol{\rho}$ 在前、旋转 $\boldsymbol{\phi}$ 在后}，与十四讲、Sophus、Barfoot 一致）。
记 $\mathrm{Exp}(\boldsymbol{\phi}):=\exp(\boldsymbol{\phi}^\wedge)$、$\mathrm{Log}(\mathbf{R}):=\ln(\mathbf{R})^\vee$（吃/吐向量的便捷算子，沿用 Handbook\cite{carlone2026handbook}、Solà\cite{sola2018micro}）。
\textbf{扰动方向：全书以右扰动 $\mathbf{R}\,\mathrm{Exp}(\delta\boldsymbol{\phi})$ 为主}（局部坐标，对接因子图/自动微分），右雅可比 $\mathbf{J}_r$ 为主；引用左扰动来源（Barfoot、十四讲）处并列给出经伴随的转换。\textbf{注意 $\boldsymbol{\rho}\neq\mathbf{t}$}（见 \cref{thm:se3-exp}）。
\end{note}

% ════════════════════════════════════════════════════════════════
\section{为什么需要李群与李代数}
\label{sec:why-lie}

\paragraph{动机：位姿是待估计量。}
在 SLAM 中，相机/机器人的位姿 $\mathbf{R},\mathbf{t}$ 并非已知，而是要从观测里\emph{估计}出来。典型做法是构造优化问题：寻找最优位姿使预测观测与实际观测的误差最小\cite{gaoxiang2019slam14}。例如已知世界系下路标 $\mathbf{p}$ 与含噪观测 $\mathbf{z}$，理想满足 $\mathbf{z}=\mathbf{T}\mathbf{p}$，于是最小化
\begin{equation}
  \min_{\mathbf{T}}\ J(\mathbf{T})=\sum_{i=1}^{N}\big\|\mathbf{z}_i-\mathbf{T}\mathbf{p}_i\big\|_2^2.
  \label{eq:pose-opt}
\end{equation}
求解 \cref{eq:pose-opt} 需要 $J$ 关于 $\mathbf{T}$ 的导数，以便迭代地调整当前估计。

\paragraph{反面：直接拿矩阵当优化变量会卡住。}
旋转矩阵自带约束：正交（$\mathbf{R}\mathbf{R}^\top=\mathbf{I}$）且 $\det\mathbf{R}=1$。若把 $\mathbf{R}$ 的 $9$ 个元素当成自由变量做梯度下降，每走一步都会\emph{破坏正交性}，迭代结果不再是旋转。强行带约束优化（如拉格朗日乘子）则代价高、易病态。问题的根子在于：$\SO(3)$、$\SE(3)$ \textbf{不是向量空间}，没有自由的加减法，求导的差商定义 $\lim(\,f(\mathbf{R}+\Delta\mathbf{R})-f(\mathbf{R})\,)/\Delta\mathbf{R}$ 里 “$\mathbf{R}+\Delta\mathbf{R}$” 根本不合法。

\paragraph{历史：从挪威数学家 Lie 到机器人学。}
$19$ 世纪挪威数学家 Sophus Lie 研究连续变换群（今称\emph{李群}）时发现：一个连续群在单位元附近的全部局部信息，可由一个\emph{线性}对象——它的\emph{李代数}——精确编码，二者经指数映射相连。机器人学自 Chirikjian\cite{chirikjian2011stochastic} 等人起系统地把这套工具引入位姿估计；今天它是 Barfoot《状态估计》\cite{barfoot2024state}、SLAM Handbook\cite{carlone2026handbook} 的共同语言，也是 GTSAM、Sophus、Ceres、manif 等库的数学内核。

\paragraph{核心思想。}
李群（一个\emph{流形}）在单位元附近的局部结构，由一个\emph{向量空间}——李代数——精确刻画；二者经指数/对数映射局部一一对应。于是：在李代数这一线性空间上自由地加减、求导，再用指数映射映回李群。这把\textbf{带约束优化转化为无约束优化}——正是全书反复使用的范式。

\begin{figure}[ht]
\centering
\begin{tikzpicture}[scale=1.5,>=Stealth]
  \draw[thick, main] (0,0) circle (1);
  \node[main] at (-1.15,0.95) {$\mathcal{M}$ (李群)};
  \coordinate (E) at (0,1);
  \fill (E) circle (0.8pt) node[above right] {$\mathcal{E}$ (单位元)};
  \draw[thick, second] (-1.1,1) -- (1.25,1) node[right] {$T_{\mathcal{E}}\mathcal{M}$ (李代数/切空间)};
  \draw[->, third, thick] (E) .. controls (0.5,0.95) and (0.85,0.6) .. (0.92,0.38);
  \node[third] at (1.2,0.55) {$\exp$};
  \fill[second] (0.6,1) circle (0.7pt) node[above] {$\boldsymbol{\phi}$};
\end{tikzpicture}
\caption{李群流形 $\mathcal{M}$ 与其在单位元 $\mathcal{E}$ 处的切空间（李代数）。切向量 $\boldsymbol{\phi}$ 经指数映射“贴合”到流形上。（仿 \cite{barfoot2024state} 重绘）}
\label{fig:lie-manifold}
\end{figure}

\begin{pitfall}[概念误区：“流形 = 弯曲，所以一定难算”]
新手常以为“活在曲面上就只能用复杂的微分几何”。实际上，流形的\emph{定义性质}恰恰是\textbf{局部像欧氏空间}（\cref{fig:lie-manifold} 的切平面）。李群之所以好用，正因为我们\emph{只在单位元附近的切空间（李代数）里做线性运算}，再映回去——难的部分被指数映射吸收了。理解这一点，本章后面所有“扰动 + 求导”的套路都顺理成章。
\end{pitfall}

% ════════════════════════════════════════════════════════════════
\section{群与矩阵李群}
\label{sec:group}

\paragraph{直觉：加法不封闭，只有乘法封闭。}
为什么旋转“只能是群、不能是向量空间”？一个朴素却关键的观察\cite{gaoxiang2019slam14}：任取两个旋转矩阵 $\mathbf{R}_1,\mathbf{R}_2$，它们的\emph{和}一般不再是旋转矩阵（例如 $\mathbf{R}+(-\mathbf{R})=\mathbf{0}\notin\SO(3)$）：
\begin{equation}
  \mathbf{R}_1+\mathbf{R}_2\notin\SO(3),\qquad \mathbf{T}_1+\mathbf{T}_2\notin\SE(3).
  \label{eq:not-vs}
\end{equation}
但它们的\emph{乘积}仍是旋转/位姿：$\mathbf{R}_1\mathbf{R}_2\in\SO(3)$，$\mathbf{T}_1\mathbf{T}_2\in\SE(3)$，对应“先后做两次变换”的复合。对于这种\emph{只配备一个良好运算（乘法）}的集合，数学上称之为\emph{群}。

\begin{definition}[群]\label{def:group}
群是集合 $A$ 与二元运算 $\cdot$ 构成的代数结构 $G=(A,\cdot)$，满足（助记“封结幺逆”）：
\begin{enumerate}
  \item \textbf{封}闭性：$\forall a_1,a_2\in A,\; a_1\cdot a_2\in A$；
  \item \textbf{结}合律：$\forall a_1,a_2,a_3\in A,\;(a_1\cdot a_2)\cdot a_3=a_1\cdot(a_2\cdot a_3)$；
  \item \textbf{幺}元：$\exists a_0\in A,\;\forall a\in A,\; a_0\cdot a=a\cdot a_0=a$；
  \item 逆：$\forall a\in A,\;\exists a^{-1}\in A,\; a\cdot a^{-1}=a_0$。
\end{enumerate}
\end{definition}

\noindent 常见群：整数加法 $(\mathbb{Z},+)$、非零有理数乘法 $(\mathbb{Q}\setminus\{0\},\cdot)$；矩阵中的一般线性群 $\mathrm{GL}(n)$、特殊正交群 $\SO(n)$、特殊欧氏群 $\SE(n)$。\textbf{李群}指\emph{具有连续（光滑）结构}的群：离散的 $(\mathbb{Z},+)$ 不是李群，而刚体可以在空间中\emph{连续}地运动，故 $\SO(n),\SE(n)$ 都是李群。三维旋转与位姿分别是：
\begin{align}
  \SO(3) &= \big\{\mathbf{R}\in\mathbb{R}^{3\times3}\,\big|\,\mathbf{R}\mathbf{R}^\top=\mathbf{I},\ \det\mathbf{R}=1\big\},\label{eq:so3}\\
  \SE(3) &= \Big\{\mathbf{T}=\begin{bmatrix}\mathbf{R}&\mathbf{t}\\\mathbf{0}^\top&1\end{bmatrix}\in\mathbb{R}^{4\times4}\,\Big|\,\mathbf{R}\in\SO(3),\ \mathbf{t}\in\mathbb{R}^3\Big\}.\label{eq:se3}
\end{align}

\begin{proposition}[$\SO(3),\SE(3)$ 是矩阵李群]\label{prop:matrix-lie}
$\SO(3),\SE(3)$ 关于矩阵乘法满足 \cref{def:group} 四条公理，且是光滑流形，故为矩阵李群；但二者对矩阵加法\emph{不封闭}（\cref{eq:not-vs}），故\emph{不是}向量空间。
\end{proposition}
\begin{proof}
结合律由矩阵乘法的结合律直接给出；幺元为 $\mathbf{I}$。\emph{封闭性}：$\SO(3)$ 中 $(\mathbf{R}_1\mathbf{R}_2)(\mathbf{R}_1\mathbf{R}_2)^\top=\mathbf{R}_1(\mathbf{R}_2\mathbf{R}_2^\top)\mathbf{R}_1^\top=\mathbf{I}$，$\det(\mathbf{R}_1\mathbf{R}_2)=\det\mathbf{R}_1\det\mathbf{R}_2=1$；$\SE(3)$ 中 $\mathbf{T}_1\mathbf{T}_2=\big[\begin{smallmatrix}\mathbf{R}_1\mathbf{R}_2&\mathbf{R}_1\mathbf{t}_2+\mathbf{t}_1\\\mathbf{0}^\top&1\end{smallmatrix}\big]$ 仍是位姿。\emph{逆元}：$\mathbf{R}^{-1}=\mathbf{R}^\top\in\SO(3)$；$\mathbf{T}^{-1}=\big[\begin{smallmatrix}\mathbf{R}^\top&-\mathbf{R}^\top\mathbf{t}\\\mathbf{0}^\top&1\end{smallmatrix}\big]$，验证 $\mathbf{T}\mathbf{T}^{-1}=\big[\begin{smallmatrix}\mathbf{R}\mathbf{R}^\top&-\mathbf{R}\mathbf{R}^\top\mathbf{t}+\mathbf{t}\\\mathbf{0}^\top&1\end{smallmatrix}\big]=\mathbf{I}$。\emph{加法不封闭}：零矩阵 $\mathbf{0}\notin\SO(3)$。
\end{proof}

\begin{note}
\textbf{$\SE(3)$ 的逆有闭式}（三书一致，常用）：
\[
  \mathbf{T}^{-1}=\begin{bmatrix}\mathbf{R}^\top&-\mathbf{R}^\top\mathbf{t}\\\mathbf{0}^\top&1\end{bmatrix},
\]
即先用 $\mathbf{R}^\top$ 反转旋转，再把平移反向并旋转回来。Barfoot 记旋转为 $\mathbf{C}$；SLAM Handbook 多从流形/因子图视角使用同一对象。
\end{note}

\begin{pitfall}[概念误区：把 $\mathbf{R}^{-1}=\mathbf{R}^\top$ 当成巧合]
“旋转矩阵的逆恰好等于转置，真方便”——这\emph{不是}巧合，而是正交矩阵的\textbf{定义性质}：$\mathbf{R}\in\SO(3)\iff\mathbf{R}^\top\mathbf{R}=\mathbf{I}$。几何上，旋转保长度（保范），所以它的逆就是“反着转”，即转置。理解了这点，就知道为什么 $\SE(3)$ 的逆\emph{不是}简单转置（平移项要额外处理）——$\SE(3)$ 不是正交群。
\end{pitfall}

\begin{practice}
\begin{enumerate}
  \item 验证 $\SO(3)$、$\SE(3)$ 关于矩阵乘法构成群（逐条核对“封结幺逆”）。
  \item 由 $\mathbf{T}=\big[\begin{smallmatrix}\mathbf{R}&\mathbf{t}\\\mathbf{0}^\top&1\end{smallmatrix}\big]$ 直接验证上面 $\mathbf{T}^{-1}$ 的闭式（提示：算 $\mathbf{T}\mathbf{T}^{-1}$）。
  \item 举一个反例说明 $\SO(3)$ 对加法不封闭。
\end{enumerate}
\end{practice}

% ════════════════════════════════════════════════════════════════
\section{李代数}
\label{sec:lie-algebra}

\paragraph{动机与招牌推导：从 $\mathbf{R}\mathbf{R}^\top=\mathbf{I}$ 自然“挤出”李代数。}
我们不直接“天降”李代数这个符号，而是顺着\emph{旋转随时间变化}把它逼出来\cite{gaoxiang2019slam14}。设旋转 $\mathbf{R}(t)$ 连续变化，恒满足约束 $\mathbf{R}(t)\mathbf{R}(t)^\top=\mathbf{I}$。两边对时间求导：
\begin{equation}
  \dot{\mathbf{R}}\mathbf{R}^\top+\mathbf{R}\dot{\mathbf{R}}^\top=\mathbf{0}
  \;\Longrightarrow\;
  \dot{\mathbf{R}}\mathbf{R}^\top=-\big(\dot{\mathbf{R}}\mathbf{R}^\top\big)^\top.
  \label{eq:RRt-deriv}
\end{equation}
\cref{eq:RRt-deriv} 说明 $\dot{\mathbf{R}}\mathbf{R}^\top$ 是\textbf{反对称矩阵}。任一反对称矩阵都对应唯一三维向量 $\boldsymbol{\phi}(t)\in\mathbb{R}^3$，记 $\dot{\mathbf{R}}\mathbf{R}^\top=\boldsymbol{\phi}^\wedge$，其中
\begin{equation}
  \boldsymbol{\phi}^\wedge=\begin{bmatrix}0&-\phi_3&\phi_2\\\phi_3&0&-\phi_1\\-\phi_2&\phi_1&0\end{bmatrix},\qquad (\boldsymbol{\phi}^\wedge)^\vee=\boldsymbol{\phi}.
  \label{eq:hat}
\end{equation}
右乘 $\mathbf{R}$（正交，$\mathbf{R}^\top\mathbf{R}=\mathbf{I}$）即得旋转的\emph{微分方程}：
\begin{equation}
  \dot{\mathbf{R}}(t)=\boldsymbol{\phi}(t)^\wedge\,\mathbf{R}(t).
  \label{eq:Rdot}
\end{equation}
直觉是：\emph{每对旋转求一次导，相当于左乘一个反对称矩阵 $\boldsymbol{\phi}^\wedge$}。在 $\mathbf{R}(0)=\mathbf{I}$ 处，$\boldsymbol{\phi}$ 刻画了 $\mathbf{R}$ 的一阶变化，正是 $\mathbf{R}$ 在单位元的\emph{切向量}。若 $\boldsymbol{\phi}$ 在该处取常值 $\boldsymbol{\phi}_0$，\cref{eq:Rdot} 是线性常微分方程，解为
\begin{equation}
  \mathbf{R}(t)=\exp(\boldsymbol{\phi}_0^\wedge\, t),
  \label{eq:exp-emerge}
\end{equation}
——旋转矩阵与反对称矩阵\emph{经矩阵指数发生了联系}。这就引出两个待解决的问题：(i) 与 $\mathbf{R}$ 对应的 $\boldsymbol{\phi}$ 是什么？它就是 $\SO(3)$ 的李代数 $\so(3)$ 元素；(ii) 矩阵指数 $\exp(\boldsymbol{\phi}^\wedge)$ 怎么算、能否反算？这就是指数/对数映射（\cref{sec:exp-log}）。

\begin{definition}[李代数]\label{def:lie-algebra}
李代数由向量空间 $\mathbb{V}$、数域 $\mathbb{F}$ 与二元运算（李括号）$[\cdot,\cdot]:\mathbb{V}\times\mathbb{V}\to\mathbb{V}$ 组成，满足：
\begin{enumerate}
  \item 封闭性：$\forall \mathbf{X},\mathbf{Y}\in\mathbb{V},\;[\mathbf{X},\mathbf{Y}]\in\mathbb{V}$；
  \item 双线性：$[a\mathbf{X}+b\mathbf{Y},\mathbf{Z}]=a[\mathbf{X},\mathbf{Z}]+b[\mathbf{Y},\mathbf{Z}]$（对第二变元同理）；
  \item 自反性：$\forall\mathbf{X}\in\mathbb{V},\;[\mathbf{X},\mathbf{X}]=\mathbf{0}$；
  \item 雅可比恒等式：$[\mathbf{X},[\mathbf{Y},\mathbf{Z}]]+[\mathbf{Z},[\mathbf{X},\mathbf{Y}]]+[\mathbf{Y},[\mathbf{Z},\mathbf{X}]]=\mathbf{0}$。
\end{enumerate}
李括号“度量两元素的差异”，不要求结合律。三维向量空间配叉积 $(\mathbb{R}^3,\mathbb{R},\times)$ 就是一个李代数。
\end{definition}

\begin{definition}[$\so(3)$ 与 $\se(3)$]\label{def:so3-se3}
\begin{align}
  \so(3) &= \big\{\boldsymbol{\phi}\in\mathbb{R}^3,\ \boldsymbol{\Phi}=\boldsymbol{\phi}^\wedge\in\mathbb{R}^{3\times3}\big\},\quad
  [\boldsymbol{\phi}_1,\boldsymbol{\phi}_2]=(\boldsymbol{\phi}_1^\wedge\boldsymbol{\phi}_2^\wedge-\boldsymbol{\phi}_2^\wedge\boldsymbol{\phi}_1^\wedge)^\vee;\\
  \se(3) &= \Big\{\boldsymbol{\xi}=\begin{bmatrix}\boldsymbol{\rho}\\\boldsymbol{\phi}\end{bmatrix}\in\mathbb{R}^6,\ \boldsymbol{\xi}^\wedge=\begin{bmatrix}\boldsymbol{\phi}^\wedge&\boldsymbol{\rho}\\\mathbf{0}^\top&0\end{bmatrix}\in\mathbb{R}^{4\times4}\Big\}.
\end{align}
$\so(3)$ 与 $\SO(3)$ 的关系由 $\mathbf{R}=\exp(\boldsymbol{\phi}^\wedge)$ 给定。$\so(3)$ 恰是“由旋转向量组成、用来表达旋转矩阵导数”的空间。
\end{definition}

\begin{note}
此处 $(\cdot)^\wedge$ 被\emph{重载}：在 $\so(3)$ 是反对称化 $\mathbb{R}^3\to\mathbb{R}^{3\times3}$；在 $\se(3)$ 是 $\mathbb{R}^6\to\mathbb{R}^{4\times4}$（右上块直接放 $\boldsymbol{\rho}$，\textbf{不再反对称}）。$\boldsymbol{\rho}$ 是 $\se(3)$ 的“平移分量”，但\textbf{并非}位姿的真实平移 $\mathbf{t}$；二者相差一个左雅可比 $\mathbf{t}=\mathbf{J}\boldsymbol{\rho}$（见 \cref{thm:se3-exp}）。
\end{note}

\begin{pitfall}[概念误区：把“李代数”当成“切空间本身”或把 $\boldsymbol{\phi}^\wedge$ 当成切空间]
两个常见混淆：(1) $\so(3)$ 的\emph{元素}是切向量 $\boldsymbol{\phi}$（或其反对称化 $\boldsymbol{\phi}^\wedge$），而\emph{整个} $\so(3)$ 才是切空间 $T_{\mathcal{E}}\SO(3)$——元素与空间别混。(2) $\boldsymbol{\phi}$ 与 $\boldsymbol{\phi}^\wedge$ 只是同一对象的“向量形式”与“矩阵形式”，靠 $\wedge/\vee$ 互换，做线性代数时常不加区分，但维度要心里有数（$3$ 维向量 vs $3\times3$ 矩阵）。
\end{pitfall}

\begin{practice}
\begin{enumerate}
  \item 验证 $(\mathbb{R}^3,\mathbb{R},\times)$ 满足 \cref{def:lie-algebra} 四条（尤其用 $\mathbf{a}\times\mathbf{a}=\mathbf{0}$ 验自反性、用叉积的雅可比恒等式验第四条）。
  \item 验证 $\so(3)$、$\se(3)$ 在各自李括号下满足李代数公理。
  \item （概念）解释为什么 \cref{eq:RRt-deriv} 一定给出反对称矩阵，而非任意矩阵。
\end{enumerate}
\end{practice}

% ════════════════════════════════════════════════════════════════
\section{指数与对数映射}
\label{sec:exp-log}

李代数到李群的桥是\emph{指数映射} $\mathbf{R}=\exp(\boldsymbol{\phi}^\wedge)$。直接按矩阵指数的级数 $\exp(\mathbf{A})=\sum_{n\ge0}\mathbf{A}^n/n!$ 计算要算无穷次幂，并不现实；幸好 $\so(3)$ 有闭式。

\begin{theorem}[$\SO(3)$ 指数映射：罗德里格斯公式]\label{thm:so3-exp}
设 $\boldsymbol{\phi}=\theta\mathbf{a}$，$\theta=\|\boldsymbol{\phi}\|$，$\mathbf{a}$ 为单位轴，则
\begin{equation}
  \mathbf{R}=\exp(\boldsymbol{\phi}^\wedge)=\cos\theta\,\mathbf{I}+(1-\cos\theta)\,\mathbf{a}\mathbf{a}^\top+\sin\theta\,\mathbf{a}^\wedge.
  \label{eq:rodrigues}
\end{equation}
\end{theorem}

\begin{proof}
关键是用两条单位向量恒等式（$\mathbf{a}$ 为单位向量）把高阶幂降阶：
\begin{equation}
  \mathbf{a}^\wedge\mathbf{a}^\wedge=\mathbf{a}\mathbf{a}^\top-\mathbf{I},\qquad \mathbf{a}^\wedge\mathbf{a}^\wedge\mathbf{a}^\wedge=-\mathbf{a}^\wedge,
  \label{eq:a-identities}
\end{equation}
于是 $\mathbf{a}^\wedge$ 的幂在 $\{\mathbf{a}^\wedge,(\mathbf{a}^\wedge)^2\}$ 间循环。代入指数级数并按奇偶项归并：
\begin{align*}
  \exp(\theta\mathbf{a}^\wedge)
  &=\mathbf{I}+\theta\mathbf{a}^\wedge+\tfrac{\theta^2}{2!}(\mathbf{a}^\wedge)^2+\tfrac{\theta^3}{3!}(\mathbf{a}^\wedge)^3+\tfrac{\theta^4}{4!}(\mathbf{a}^\wedge)^4+\cdots\\
  &=\mathbf{a}\mathbf{a}^\top+\Big(\theta-\tfrac{\theta^3}{3!}+\tfrac{\theta^5}{5!}-\cdots\Big)\mathbf{a}^\wedge
   -\Big(1-\tfrac{\theta^2}{2!}+\tfrac{\theta^4}{4!}-\cdots\Big)(\mathbf{a}^\wedge)^2\\
  &=\mathbf{a}\mathbf{a}^\top+\sin\theta\,\mathbf{a}^\wedge-\cos\theta\,(\mathbf{a}\mathbf{a}^\top-\mathbf{I})
   =\cos\theta\,\mathbf{I}+(1-\cos\theta)\,\mathbf{a}\mathbf{a}^\top+\sin\theta\,\mathbf{a}^\wedge,
\end{align*}
末步用 $(\mathbf{a}^\wedge)^2=\mathbf{a}\mathbf{a}^\top-\mathbf{I}$ 回代，即 \cref{eq:rodrigues}。
\end{proof}
\noindent 注意 \cref{eq:rodrigues} 与 \cref{ch:rigid_body} 的罗德里格斯公式\emph{一模一样}——这正说明 $\so(3)$ 就是旋转向量的空间，指数映射就是罗德里格斯公式。上面的证明是\emph{代数}的（级数按奇偶项归并）；其\emph{几何}诠释见 \cref{fig:rodrigues}。

\begin{figure}[ht]
  \centering
  \includegraphics[width=.52\linewidth]{figures/pdf/rodrigues.pdf}
  \caption{罗德里格斯公式的几何诠释：向量 $\mathbf{v}$ 绕单位轴 $\mathbf{k}$ 转 $\theta$。把 $\mathbf{v}$ 分解为平行分量 $\mathbf{v}_\parallel$（沿轴、旋转不变）与垂直分量 $\mathbf{v}_\perp$；后者在垂直平面内转 $\theta$ 得 $\mathbf{v}_\perp'=\cos\theta\,\mathbf{v}_\perp+\sin\theta\,(\mathbf{k}\times\mathbf{v}_\perp)$，再与 $\mathbf{v}_\parallel$ 叠加即得 $\mathbf{v}_{\mathrm{rot}}$。这正是 \cref{eq:rodrigues} 的向量（几何）形式。}
  \label{fig:rodrigues}
\end{figure}

\paragraph{对数映射（反向）。}由 $\mathbf{R}$ 反解 $\boldsymbol{\phi}=\mathrm{Log}(\mathbf{R})$，不必用级数，用迹与特征向量更省：
\begin{equation}
  \theta=\cos^{-1}\!\frac{\operatorname{tr}(\mathbf{R})-1}{2},\qquad \mathbf{R}\mathbf{a}=\mathbf{a}\ \ (\mathbf{a}\text{ 为特征值 }1\text{ 的单位特征向量}).
  \label{eq:so3-log}
\end{equation}

\begin{proposition}[满射非单射]\label{prop:surjective}
指数映射 $\so(3)\to\SO(3)$ 是\emph{满射但非单射}：$\theta$ 与 $\theta+2\pi m$ 给出同一 $\mathbf{R}$。若限制 $|\theta|<\pi$，则李群与李代数元素一一对应。
\end{proposition}

\begin{pitfall}[思维陷阱：忽视周期性，优化跨越 $\pm\pi$ 处“跳变”]
因为指数映射非单射（\cref{prop:surjective}），$\mathrm{Log}$ 通常取主值 $|\theta|\le\pi$。当真实旋转接近 $\pi$、或优化迭代跨过该边界时，$\mathrm{Log}$ 会发生 $2\pi$ 跳变，残差和雅可比出现不连续。实现里常见后果是 BA/位姿图在大旋转处不收敛。对策：增量更新（扰动）每步都是小量，远离 $\pi$；必要时对相对旋转取 $\mathrm{Log}$ 而非绝对旋转。
\end{pitfall}

\begin{theorem}[$\SE(3)$ 指数映射]\label{thm:se3-exp}
$\mathbf{T}=\exp(\boldsymbol{\xi}^\wedge)=\big[\begin{smallmatrix}\mathbf{R}&\mathbf{J}\boldsymbol{\rho}\\\mathbf{0}^\top&1\end{smallmatrix}\big]$，其中 $\mathbf{R}=\exp(\boldsymbol{\phi}^\wedge)$，\textbf{真实平移} $\mathbf{t}=\mathbf{J}\boldsymbol{\rho}$，左雅可比
\begin{equation}
  \mathbf{J}=\sum_{n=0}^{\infty}\frac{1}{(n+1)!}(\boldsymbol{\phi}^\wedge)^n
  =\frac{\sin\theta}{\theta}\mathbf{I}+\Big(1-\frac{\sin\theta}{\theta}\Big)\mathbf{a}\mathbf{a}^\top+\frac{1-\cos\theta}{\theta}\mathbf{a}^\wedge.
  \label{eq:left-jacobian}
\end{equation}
\end{theorem}
\begin{proof}
$\exp(\boldsymbol{\xi}^\wedge)$ 按 $4\times4$ 分块展开：左上块是 $\mathbf{R}=\exp(\boldsymbol{\phi}^\wedge)$，右上块为 $\big(\sum_{n\ge0}\tfrac{1}{(n+1)!}(\boldsymbol{\phi}^\wedge)^n\big)\boldsymbol{\rho}=\mathbf{J}\boldsymbol{\rho}$。令 $\boldsymbol{\phi}=\theta\mathbf{a}$，用 $(\mathbf{a}^\wedge)^2=\mathbf{a}\mathbf{a}^\top-\mathbf{I}$、$(\mathbf{a}^\wedge)^3=-\mathbf{a}^\wedge$ 归并：
\begin{align*}
  \mathbf{J}&=\mathbf{I}+\tfrac{1}{2!}\theta\mathbf{a}^\wedge+\tfrac{1}{3!}\theta^2(\mathbf{a}^\wedge)^2+\tfrac{1}{4!}\theta^3(\mathbf{a}^\wedge)^3+\cdots\\
  &=\mathbf{I}+\frac{1-\cos\theta}{\theta}\mathbf{a}^\wedge+\frac{\theta-\sin\theta}{\theta}(\mathbf{a}\mathbf{a}^\top-\mathbf{I})
   =\frac{\sin\theta}{\theta}\mathbf{I}+\Big(1-\frac{\sin\theta}{\theta}\Big)\mathbf{a}\mathbf{a}^\top+\frac{1-\cos\theta}{\theta}\mathbf{a}^\wedge,
\end{align*}
即 \cref{eq:left-jacobian}。其逆见 \cref{eq:Jl}，在 $\theta=2\pi m\,(m\neq0)$ 处奇异。
\end{proof}

\begin{insight}
$\mathbf{t}=\mathbf{J}\boldsymbol{\rho}$ 正是“$\se(3)$ 平移分量 $\boldsymbol{\rho}\neq$ 真实平移 $\mathbf{t}$”的根源：只有当 $\boldsymbol{\phi}\to\mathbf{0}$（即 $\mathbf{J}\to\mathbf{I}$）时二者才相等。这个 $\mathbf{J}$ 不是巧合地“顺便出现”，它就是下一节 BCH 里的\emph{左雅可比}，会在扰动、伴随、不确定度里反复登场——记住它。
\end{insight}

\begin{practice}
\begin{enumerate}
  \item 用 \cref{eq:a-identities} 验证 $(\mathbf{a}^\wedge)^4=-(\mathbf{a}^\wedge)^2$，并据此说明为何级数只剩 $\sin/\cos$ 两族。
  \item （推导，在草稿纸上完成）从级数出发独立推出 \cref{eq:rodrigues}，对照正文证明。
  \item 取 $\theta\to0$ 对 \cref{eq:left-jacobian} 做泰勒展开，验证 $\mathbf{J}\to\mathbf{I}+\tfrac12\boldsymbol{\phi}^\wedge+O(\theta^2)$，从而 $\mathbf{t}\approx\boldsymbol{\rho}$。
\end{enumerate}
\end{practice}

% ════════════════════════════════════════════════════════════════
\section{BCH 公式与左右雅可比}
\label{sec:bch}

\paragraph{动机：李代数上的“加法”对应李群上的什么？}
优化要在李代数里做加减，但李群里是\emph{乘法}。于是要问：$\exp(\boldsymbol{\phi}_1^\wedge)\exp(\boldsymbol{\phi}_2^\wedge)$ 是否等于 $\exp((\boldsymbol{\phi}_1+\boldsymbol{\phi}_2)^\wedge)$？对\emph{标量}成立，但对\emph{矩阵}一般\textbf{不成立}（因为矩阵不交换）。差多少，由 Baker–Campbell–Hausdorff（BCH）公式给出：
\begin{equation}
  \ln\big(\exp(\mathbf{A})\exp(\mathbf{B})\big)=\mathbf{A}+\mathbf{B}+\tfrac12[\mathbf{A},\mathbf{B}]+\tfrac1{12}[\mathbf{A},[\mathbf{A},\mathbf{B}]]-\tfrac1{12}[\mathbf{B},[\mathbf{A},\mathbf{B}]]+\cdots
  \label{eq:bch}
\end{equation}
即乘积的对数 = 两者之和 + 一串\emph{李括号}余项。

\paragraph{线性近似——本节真正要用的。}
当其中一个为小量时，二阶以上李括号可略，BCH 退化为线性形式，并\emph{自然分出左、右两种}\cite{gaoxiang2019slam14}：
\begin{equation}
  \ln\!\big(\exp(\boldsymbol{\phi}_1^\wedge)\exp(\boldsymbol{\phi}_2^\wedge)\big)^\vee\approx
  \begin{cases}
    \mathbf{J}_l(\boldsymbol{\phi}_2)^{-1}\,\boldsymbol{\phi}_1+\boldsymbol{\phi}_2, & \boldsymbol{\phi}_1\to\mathbf{0}\ (\text{左乘小量}),\\[2pt]
    \mathbf{J}_r(\boldsymbol{\phi}_1)^{-1}\,\boldsymbol{\phi}_2+\boldsymbol{\phi}_1, & \boldsymbol{\phi}_2\to\mathbf{0}\ (\text{右乘小量}).
  \end{cases}
  \label{eq:bch-linear}
\end{equation}
解读：给 $\mathbf{R}_2$（李代数 $\boldsymbol{\phi}_2$）\emph{左乘}一个微小旋转（李代数 $\boldsymbol{\phi}_1$），在李代数上约等于在 $\boldsymbol{\phi}_2$ 上加 $\mathbf{J}_l^{-1}\boldsymbol{\phi}_1$；\emph{右乘}小量则加 $\mathbf{J}_r^{-1}$ 项。\textbf{用左乘模型还是右乘模型，须自始至终一致}——这正对应本书选定的\emph{右扰动}（\cref{sec:perturb}）。

\paragraph{左/右雅可比的闭式。}左雅可比 $\mathbf{J}_l$ 即 \cref{eq:left-jacobian} 的 $\mathbf{J}$；其逆与右版本为（$\boldsymbol{\phi}=\theta\mathbf{a}$）：
\begin{align}
  \mathbf{J}_l &= \frac{\sin\theta}{\theta}\mathbf{I}+\Big(1-\frac{\sin\theta}{\theta}\Big)\mathbf{a}\mathbf{a}^\top+\frac{1-\cos\theta}{\theta}\mathbf{a}^\wedge,\label{eq:Jl}\\
  \mathbf{J}_l^{-1} &= \frac{\theta}{2}\cot\frac{\theta}{2}\,\mathbf{I}+\Big(1-\frac{\theta}{2}\cot\frac{\theta}{2}\Big)\mathbf{a}\mathbf{a}^\top-\frac{\theta}{2}\mathbf{a}^\wedge,\\
  \mathbf{J}_r(\boldsymbol{\phi}) &= \mathbf{J}_l(-\boldsymbol{\phi})=\frac{\sin\theta}{\theta}\mathbf{I}+\Big(1-\frac{\sin\theta}{\theta}\Big)\mathbf{a}\mathbf{a}^\top-\frac{1-\cos\theta}{\theta}\mathbf{a}^\wedge,\label{eq:Jr}\\
  \mathbf{J}_r^{-1} &= \frac{\theta}{2}\cot\frac{\theta}{2}\,\mathbf{I}+\Big(1-\frac{\theta}{2}\cot\frac{\theta}{2}\Big)\mathbf{a}\mathbf{a}^\top+\frac{\theta}{2}\mathbf{a}^\wedge.
\end{align}
关键关系（贯穿全章）：$\boxed{\,\mathbf{J}_r(\boldsymbol{\phi})=\mathbf{J}_l(-\boldsymbol{\phi})=\mathbf{J}_l(\boldsymbol{\phi})^\top\,}$（最后一个等号因 $\mathbf{a}\mathbf{a}^\top$、$\mathbf{I}$ 对称而 $\mathbf{a}^\wedge$ 反对称）。当 $\theta\to0$，$\mathbf{J}_l,\mathbf{J}_r\to\mathbf{I}$，左右扰动趋于一致。

\paragraph{$\SE(3)$ 的 $6\times6$ 雅可比（Barfoot 闭式，独立记录）。}
十四讲在此略去了 $\SE(3)$ 的 $6\times6$ 雅可比；为使本书\emph{自包含}，这里补上 Barfoot\cite{barfoot2024state} 的闭式。记 $\boldsymbol{\mathcal{J}}_l(\boldsymbol{\xi})=\sum_{n\ge0}\frac{1}{(n+1)!}(\boldsymbol{\xi}^\curlywedge)^n$（$\boldsymbol{\xi}^\curlywedge$ 见 \cref{sec:adjoint}），其分块为
\begin{equation}
  \boldsymbol{\mathcal{J}}_l(\boldsymbol{\xi})=\begin{bmatrix}\mathbf{J}_l&\mathbf{Q}_l\\\mathbf{0}&\mathbf{J}_l\end{bmatrix},\qquad
  \boldsymbol{\mathcal{J}}_l^{-1}=\begin{bmatrix}\mathbf{J}_l^{-1}&-\mathbf{J}_l^{-1}\mathbf{Q}_l\mathbf{J}_l^{-1}\\\mathbf{0}&\mathbf{J}_l^{-1}\end{bmatrix},
  \label{eq:se3-6x6}
\end{equation}
其中 $\mathbf{J}_l$ 是 \cref{eq:Jl} 的 $3\times3$ 左雅可比，$\mathbf{Q}_l(\boldsymbol{\rho},\boldsymbol{\phi})$ 的完整闭式（Barfoot 式 8.91）见 \cref{sec:lie-appendix}。右版本 $\boldsymbol{\mathcal{J}}_r(\boldsymbol{\xi})=\boldsymbol{\mathcal{J}}_l(-\boldsymbol{\xi})$。一般 SLAM 旋转较小时常用近似 $\boldsymbol{\mathcal{J}}\approx\mathbf{I}$；需要高精度协方差传播时才用完整式。

\begin{pitfall}[思维陷阱：把 BCH 当成精确等式滥用]
\cref{eq:bch-linear} 是\emph{一阶近似}，只在扰动\emph{小}时成立。有人把它当精确公式用于大角度，会引入系统误差。正确用法：扰动量每步都保持小（这正是迭代优化的天然条件）；要精确合并两个有限旋转，请用群乘法 $\mathrm{Exp}(\boldsymbol{\phi}_1)\mathrm{Exp}(\boldsymbol{\phi}_2)$ 本身，或伴随恒等式（\cref{sec:adjoint}），而非 BCH 截断。
\end{pitfall}

\begin{practice}
\begin{enumerate}
  \item 由 $\mathbf{J}_r(\boldsymbol{\phi})=\mathbf{J}_l(-\boldsymbol{\phi})$ 推出 \cref{eq:Jr}，并验证 $\mathbf{J}_r=\mathbf{J}_l^\top$。
  \item （开放思考）当 $\theta\to0$ 时 $\tfrac{\theta}{2}\cot\tfrac{\theta}{2}\to1$，验证 $\mathbf{J}_l^{-1}\to\mathbf{I}-\tfrac12\boldsymbol{\phi}^\wedge$；解释这与“小扰动下左右一致”如何吻合。
\end{enumerate}
\end{practice}

% ════════════════════════════════════════════════════════════════
\section{扰动模型（右扰动为主）}
\label{sec:perturb}

要解 \cref{eq:pose-opt}，须计算误差对位姿的导数，例如旋转后点 $\mathbf{R}\mathbf{p}$ 对旋转的导数。$\SO(3)$ 无加法，差商定义失效，有两条出路\cite{gaoxiang2019slam14}：

\begin{enumerate}
  \item \textbf{李代数求导（导数模型）}：用 $\boldsymbol{\phi}$ 表示 $\mathbf{R}=\exp(\boldsymbol{\phi}^\wedge)$，对 $\boldsymbol{\phi}$ 求导。
  \item \textbf{扰动模型}：对 $\mathbf{R}$ 左乘或右乘一个微小扰动 $\mathrm{Exp}(\delta\boldsymbol{\phi})$，对扰动 $\delta\boldsymbol{\phi}$ 求导。
\end{enumerate}

\paragraph{导数模型为何不实用。}对 $\partial(\exp(\boldsymbol{\phi}^\wedge)\mathbf{p})/\partial\boldsymbol{\phi}$，因 $\SO(3)$ 无加法，对李代数加扰动 $\delta\boldsymbol{\phi}$ 并用 BCH 的左乘形式 $\exp((\boldsymbol{\phi}+\delta\boldsymbol{\phi})^\wedge)\approx\exp((\mathbf{J}_l\delta\boldsymbol{\phi})^\wedge)\exp(\boldsymbol{\phi}^\wedge)$：
\begin{align}
  \frac{\partial(\mathbf{R}\mathbf{p})}{\partial\boldsymbol{\phi}}
  &=\lim_{\delta\boldsymbol{\phi}\to\mathbf{0}}\frac{\exp((\boldsymbol{\phi}+\delta\boldsymbol{\phi})^\wedge)\mathbf{p}-\exp(\boldsymbol{\phi}^\wedge)\mathbf{p}}{\delta\boldsymbol{\phi}}
   =\lim_{\delta\boldsymbol{\phi}\to\mathbf{0}}\frac{\big(\mathbf{I}+(\mathbf{J}_l\delta\boldsymbol{\phi})^\wedge\big)\mathbf{R}\mathbf{p}-\mathbf{R}\mathbf{p}}{\delta\boldsymbol{\phi}}\notag\\
  &=\lim_{\delta\boldsymbol{\phi}\to\mathbf{0}}\frac{(\mathbf{J}_l\delta\boldsymbol{\phi})^\wedge\mathbf{R}\mathbf{p}}{\delta\boldsymbol{\phi}}
   =\lim_{\delta\boldsymbol{\phi}\to\mathbf{0}}\frac{-(\mathbf{R}\mathbf{p})^\wedge\mathbf{J}_l\delta\boldsymbol{\phi}}{\delta\boldsymbol{\phi}}
   =-(\mathbf{R}\mathbf{p})^\wedge\,\mathbf{J}_l,
  \label{eq:deriv-model}
\end{align}
第三步用 $(\mathbf{J}_l\delta\boldsymbol{\phi})^\wedge\mathbf{R}\mathbf{p}=-(\mathbf{R}\mathbf{p})^\wedge\mathbf{J}_l\delta\boldsymbol{\phi}$（叉积反交换）。结果里\emph{残留一个左雅可比 $\mathbf{J}_l$}——形式复杂、还要算 $\mathbf{J}_l$，不划算。扰动模型则把它消掉。

\paragraph{右扰动模型（本书主线）。}给 $\mathbf{R}$ \textbf{右乘}小扰动 $\mathrm{Exp}(\boldsymbol{\varphi})$（$\boldsymbol{\varphi}\to\mathbf{0}$，李代数记 $\boldsymbol{\varphi}$ 以区别于位姿的 $\boldsymbol{\phi}$），对 $\boldsymbol{\varphi}$ 求导：
\begin{align}
  \frac{\partial(\mathbf{R}\mathbf{p})}{\partial\boldsymbol{\varphi}}
  &=\lim_{\boldsymbol{\varphi}\to\mathbf{0}}\frac{\mathbf{R}\,\mathrm{Exp}(\boldsymbol{\varphi})\,\mathbf{p}-\mathbf{R}\mathbf{p}}{\boldsymbol{\varphi}}
   =\lim_{\boldsymbol{\varphi}\to\mathbf{0}}\frac{\mathbf{R}(\mathbf{I}+\boldsymbol{\varphi}^\wedge)\mathbf{p}-\mathbf{R}\mathbf{p}}{\boldsymbol{\varphi}}\notag\\
  &=\lim_{\boldsymbol{\varphi}\to\mathbf{0}}\frac{\mathbf{R}\,\boldsymbol{\varphi}^\wedge\mathbf{p}}{\boldsymbol{\varphi}}
   =\lim_{\boldsymbol{\varphi}\to\mathbf{0}}\frac{-\mathbf{R}\,\mathbf{p}^\wedge\boldsymbol{\varphi}}{\boldsymbol{\varphi}}
   =\boxed{-\mathbf{R}\,\mathbf{p}^\wedge}.
  \label{eq:right-pert-so3}
\end{align}
其中第三步用一阶泰勒 $\mathrm{Exp}(\boldsymbol{\varphi})\approx\mathbf{I}+\boldsymbol{\varphi}^\wedge$，第四步用 $\boldsymbol{\varphi}^\wedge\mathbf{p}=-\mathbf{p}^\wedge\boldsymbol{\varphi}$（叉积反交换）。结果 \cref{eq:right-pert-so3} \textbf{不含 $\mathbf{J}_l$}，干净——这就是扰动模型的优势。

\paragraph{$\SE(3)$ 右扰动。}对位姿作用齐次点 $\tilde{\mathbf{p}}=[\mathbf{p}^\top,1]^\top$，给 $\mathbf{T}$ 右乘 $\mathrm{Exp}(\delta\boldsymbol{\xi})$（$\delta\boldsymbol{\xi}=[\delta\boldsymbol{\rho};\delta\boldsymbol{\phi}]$）：
\begin{equation}
  \frac{\partial(\mathbf{T}\tilde{\mathbf{p}})}{\partial\delta\boldsymbol{\xi}}
  =\mathbf{T}\,\tilde{\mathbf{p}}^{\,\odot}
  =\begin{bmatrix}\mathbf{R}&-\mathbf{R}\,\mathbf{p}^\wedge\\\mathbf{0}^\top&\mathbf{0}^\top\end{bmatrix},
  \qquad
  \tilde{\mathbf{p}}^{\,\odot}=\begin{bmatrix}\mathbf{p}\\1\end{bmatrix}^{\!\odot}=\begin{bmatrix}\mathbf{I}&-\mathbf{p}^\wedge\\\mathbf{0}^\top&\mathbf{0}^\top\end{bmatrix},
  \label{eq:right-pert-se3}
\end{equation}
其中 $(\cdot)^\odot$ 是把齐次点变成 $4\times6$ 矩阵的算子（Handbook\cite{carlone2026handbook} 记号），满足恒等式 $\boldsymbol{\xi}^\wedge\tilde{\mathbf{p}}=\tilde{\mathbf{p}}^{\,\odot}\boldsymbol{\xi}$。推导与 \cref{eq:right-pert-so3} 同理：$\mathbf{T}\,\mathrm{Exp}(\delta\boldsymbol{\xi})\tilde{\mathbf{p}}\approx\mathbf{T}(\mathbf{I}+\delta\boldsymbol{\xi}^\wedge)\tilde{\mathbf{p}}=\mathbf{T}\tilde{\mathbf{p}}+\mathbf{T}\,\delta\boldsymbol{\xi}^\wedge\tilde{\mathbf{p}}$，故增量为 $\mathbf{T}\,\delta\boldsymbol{\xi}^\wedge\tilde{\mathbf{p}}=\mathbf{T}\,\tilde{\mathbf{p}}^{\,\odot}\delta\boldsymbol{\xi}$，对 $\delta\boldsymbol{\xi}$ 求导即得 \cref{eq:right-pert-se3}（其中 $\tilde{\mathbf{p}}^{\,\odot}$ 的右上块 $-\mathbf{p}^\wedge$ 来自 $\boldsymbol{\phi}^\wedge\mathbf{p}=-\mathbf{p}^\wedge\boldsymbol{\phi}$，左块 $\mathbf{I}$ 来自 $\boldsymbol{\rho}$）。

\paragraph{与左扰动的对照。}十四讲、Barfoot 多用\emph{左}扰动，其结果为 $\partial(\mathbf{R}\mathbf{p})/\partial\boldsymbol{\varphi}=-(\mathbf{R}\mathbf{p})^\wedge$、$\partial(\mathbf{T}\tilde{\mathbf{p}})/\partial\delta\boldsymbol{\xi}=(\mathbf{T}\tilde{\mathbf{p}})^\odot=\big[\begin{smallmatrix}\mathbf{I}&-(\mathbf{R}\mathbf{p}+\mathbf{t})^\wedge\\\mathbf{0}^\top&\mathbf{0}^\top\end{smallmatrix}\big]$。左、右结果相差一个 $\mathbf{R}$（伴随）的位置，可用 $(\mathbf{R}\mathbf{p})^\wedge=\mathbf{R}\mathbf{p}^\wedge\mathbf{R}^\top$ 互验：$-(\mathbf{R}\mathbf{p})^\wedge=-\mathbf{R}\mathbf{p}^\wedge\mathbf{R}^\top=(-\mathbf{R}\mathbf{p}^\wedge)\mathbf{R}^\top$。二者通过伴随 $\mathrm{Ad}(\mathbf{R})=\mathbf{R}$ 精确转换（\cref{sec:adjoint}）。

\begin{insight}[为什么右扰动 = 在“机体/局部系”下扰动]
Solà\cite{sola2018micro} 给了一个极清晰的图像：写成 $\mathbf{R}\,\mathrm{Exp}(\boldsymbol{\varphi})$（扰动在\emph{右}），增量 $\boldsymbol{\varphi}$ 表达在 $\mathbf{R}$ \emph{所在的局部坐标系}（机体系）；写成 $\mathrm{Exp}(\boldsymbol{\varphi})\,\mathbf{R}$（扰动在\emph{左}），$\boldsymbol{\varphi}$ 表达在\emph{全局/世界系}（单位元处）。SLAM 里 IMU 角速度、相机增量观测天然在机体系给出，故\emph{右扰动}与传感器/求解器（GTSAM、Ceres、Sophus 的右乘更新）天然对齐——这就是本书选右扰动为主的根本理由。
\end{insight}

\paragraph{优化中的位姿更新。}于是迭代里位姿的“无约束更新”写成右乘指数：
\begin{equation}
  \mathbf{R}\leftarrow\mathbf{R}\,\mathrm{Exp}(\delta\boldsymbol{\phi}^*),\qquad
  \mathbf{T}\leftarrow\mathbf{T}\,\mathrm{Exp}(\delta\boldsymbol{\xi}^*),
  \label{eq:right-update}
\end{equation}
其中 $\delta\boldsymbol{\phi}^*,\delta\boldsymbol{\xi}^*$ 是在切空间（李代数）解出的增量。这保证更新后仍严格在 $\SO(3)/\SE(3)$ 上——把带约束优化变成了切空间里的\emph{无约束}最小二乘。

\begin{pitfall}[编程陷阱：左右扰动混用导致雅可比与更新不一致]
最常见的实现 bug：雅可比按\emph{右}扰动推（\cref{eq:right-pert-so3}），更新却写成\emph{左}乘 $\mathrm{Exp}(\delta)\mathbf{R}$（或反之）。后果是收敛极慢甚至发散，且难定位。铁律：\textbf{雅可比的扰动侧、迭代更新的乘法侧、协方差定义的扰动侧，三者必须同侧}。本书统一右侧；引用左扰动公式时务必先用伴随转成右侧再用。
\end{pitfall}

\begin{practice}
\begin{enumerate}
  \item （推导，在草稿纸上完成）仿照 \cref{eq:right-pert-so3}，推导\emph{左}扰动下 $\partial(\mathbf{R}\mathbf{p})/\partial\boldsymbol{\varphi}=-(\mathbf{R}\mathbf{p})^\wedge$，并用 $(\mathbf{R}\mathbf{p})^\wedge=\mathbf{R}\mathbf{p}^\wedge\mathbf{R}^\top$ 与右扰动结果对照。
  \item 推导 \cref{eq:right-pert-se3}（提示：$\mathrm{Exp}(\delta\boldsymbol{\xi})\approx\mathbf{I}+\delta\boldsymbol{\xi}^\wedge$，按齐次坐标分块，用 $\boldsymbol{\xi}^\wedge\tilde{\mathbf{p}}=\tilde{\mathbf{p}}^{\,\odot}\boldsymbol{\xi}$）。
  \item （开放思考）若观测误差定义为 $\mathbf{e}=\mathrm{Log}(\mathbf{T}^{-1}\mathbf{Z})$（位姿测量 $\mathbf{Z}$），用右扰动 $\mathbf{T}\leftarrow\mathbf{T}\,\mathrm{Exp}(\delta\boldsymbol{\xi})$ 推 $\partial\mathbf{e}/\partial\delta\boldsymbol{\xi}$（结果含 $\boldsymbol{\mathcal{J}}_r^{-1}$）。
\end{enumerate}
\end{practice}

% ════════════════════════════════════════════════════════════════
\section{伴随：左右扰动的精确桥梁}
\label{sec:adjoint}

伴随回答一个反复出现的问题：把一个李代数元素“搬过”一个群元素，会变成什么？它是\emph{左右扰动转换、坐标系变换、协方差传播}的统一工具。

\paragraph{群伴随 $\mathrm{Ad}$。}定义 $\mathrm{Ad}(\mathbf{T})$ 使共轭 $\mathbf{T}\,\boldsymbol{\xi}^\wedge\mathbf{T}^{-1}=(\mathrm{Ad}(\mathbf{T})\boldsymbol{\xi})^\wedge$。其闭式为（Barfoot\cite{barfoot2024state}、Handbook\cite{carlone2026handbook}）
\begin{equation}
  \mathrm{Ad}(\mathbf{R})=\mathbf{R}\ \ (\SO(3)),\qquad
  \mathrm{Ad}(\mathbf{T})=\begin{bmatrix}\mathbf{R}&\mathbf{t}^\wedge\mathbf{R}\\\mathbf{0}&\mathbf{R}\end{bmatrix}\in\mathbb{R}^{6\times6}\ \ (\SE(3)).
  \label{eq:Ad}
\end{equation}
最有用的是\textbf{精确}（无近似）的“移动指数”恒等式：
\begin{equation}
  \mathbf{T}\,\mathrm{Exp}(\boldsymbol{\xi})=\mathrm{Exp}\big(\mathrm{Ad}(\mathbf{T})\boldsymbol{\xi}\big)\,\mathbf{T}.
  \label{eq:Ad-move}
\end{equation}
\cref{eq:Ad-move} 把\emph{右}扰动 $\boldsymbol{\xi}$ 精确换成\emph{左}扰动 $\mathrm{Ad}(\mathbf{T})\boldsymbol{\xi}$（反之亦然）。这正是“引用左扰动公式后转成右扰动”的依据，也是把扰动跨越已知变换搬移的标准手段。

\paragraph{小伴随 $\mathrm{ad}$（$\boldsymbol{\xi}^\curlywedge$）。}群伴随在单位元的导数即\emph{小伴随}，等价于李括号 $\mathrm{ad}(\boldsymbol{\xi}^\wedge)\mathbf{x}^\wedge=[\boldsymbol{\xi}^\wedge,\mathbf{x}^\wedge]$，记输出向量为 $\boldsymbol{\xi}^\curlywedge\mathbf{x}$：
\begin{equation}
  \boldsymbol{\phi}^\curlywedge=\boldsymbol{\phi}^\wedge\ \ (\so(3)),\qquad
  \boldsymbol{\xi}^\curlywedge=\begin{bmatrix}\boldsymbol{\rho}\\\boldsymbol{\phi}\end{bmatrix}^{\!\curlywedge}=\begin{bmatrix}\boldsymbol{\phi}^\wedge&\boldsymbol{\rho}^\wedge\\\mathbf{0}&\boldsymbol{\phi}^\wedge\end{bmatrix}\in\mathbb{R}^{6\times6}\ \ (\se(3)).
  \label{eq:ad}
\end{equation}
两者经指数相连：$\mathrm{Ad}(\mathrm{Exp}(\boldsymbol{\xi}))=\exp(\mathrm{ad}(\boldsymbol{\xi}))$，即 $\mathrm{Ad}(\mathbf{T})=\exp(\boldsymbol{\xi}^\curlywedge)$。这也解释了上节 $6\times6$ 雅可比 \cref{eq:se3-6x6} 里 $\boldsymbol{\xi}^\curlywedge$ 的来历。

\begin{pitfall}[概念误区：以为 $\mathrm{Ad}(\mathbf{T})$ 也是个“变换矩阵”]
$\mathrm{Ad}(\mathbf{T})$ 是 $6\times6$、作用在\emph{李代数向量 $\boldsymbol{\xi}\in\mathbb{R}^6$}上的线性映射，不是作用在点上的 $4\times4$ 位姿。它把“切空间里的扰动”从一个坐标系搬到另一个。混淆二者会在协方差传播（\cref{sec:uncertainty}）里把维度和对象用错。
\end{pitfall}

\begin{practice}
\begin{enumerate}
  \item 证明 $\SO(3)$ 伴随性质 $\mathbf{R}\,\mathrm{Exp}(\boldsymbol{\phi})\,\mathbf{R}^\top=\mathrm{Exp}(\mathbf{R}\boldsymbol{\phi})$，即 $\mathbf{R}\boldsymbol{\phi}^\wedge\mathbf{R}^\top=(\mathbf{R}\boldsymbol{\phi})^\wedge$。
  \item 由 \cref{eq:Ad-move} 推出：右扰动雅可比 $\mathbf{J}^{\text{右}}$ 与左扰动雅可比 $\mathbf{J}^{\text{左}}$ 满足 $\mathbf{J}^{\text{左}}=\mathbf{J}^{\text{右}}\,\mathrm{Ad}(\mathbf{T})^{-1}$（对 $\mathbf{T}\tilde{\mathbf{p}}$ 的两种结果做对照）。
\end{enumerate}
\end{practice}

% ════════════════════════════════════════════════════════════════
\section{流形视角：retraction 与 \texorpdfstring{$\boxplus\,\boxminus$}{boxplus/boxminus}（对接因子图）}
\label{sec:manifold}

前面用“群 + 李代数 + 指数”讲清了机理。现代 SLAM（因子图、GTSAM、Ceres）用一套更抽象、更可复用的语言把这一切统一起来——这套语言也解释了你在求解器里必须填的“流形局部参数化”到底是什么\cite{carlone2026handbook,hertzberg2013integrating}。

\paragraph{流形、切空间、局部坐标。}状态变量带约束（如 $\mathbf{R}\mathbf{R}^\top=\mathbf{I}$），约束把变量限制在一个\emph{流形} $\mathcal{M}$ 上。流形在每点 $\boldsymbol{\chi}$ 附近\emph{局部像欧氏空间}：其\emph{切空间} $T_{\boldsymbol{\chi}}\mathcal{M}$ 配一组基，给出该点的\emph{局部坐标}。优化时，我们不动 $\mathcal{M}$ 上的点，而是在切空间里解一个普通的欧氏增量 $\boldsymbol{\xi}$，再“贴回”流形。

\paragraph{retraction（收缩映射）。}把切空间增量映回流形的映射称为 retraction。指数映射 $\mathrm{Exp}$ 是\emph{首选}的 retraction，但\textbf{并非唯一}——任何满足一阶相切条件的映射都可用（Handbook、Boumal）。这把李群优化纳入更一般的\emph{黎曼优化}框架。

\paragraph{$\boxplus/\boxminus$ 算子（求解器接口）。}定义“李群元素 $\boxplus$ 切向量”与“两元素相 $\boxminus$ 得切向量”（本书取右乘约定，与 GTSAM 默认一致；Handbook 记作 $\oplus/\ominus$）：
\begin{equation}
  \mathbf{T}\boxplus\boldsymbol{\xi}:=\mathbf{T}\,\mathrm{Exp}(\boldsymbol{\xi}),\qquad
  \mathbf{T}_2\boxminus\mathbf{T}_1:=\mathrm{Log}\big(\mathbf{T}_1^{-1}\mathbf{T}_2\big),
  \label{eq:boxplus}
\end{equation}
满足互逆 $(\mathbf{T}\boxplus\boldsymbol{\xi})\boxminus\mathbf{T}=\boldsymbol{\xi}$。有了 $\boxplus/\boxminus$，\emph{流形优化在形式上与欧氏优化无异}：误差线性化只剩对切向量 $\boldsymbol{\xi}$ 的标准雅可比，解出 $\boldsymbol{\xi}^*$ 后用 $\boxplus$ 收缩回流形（\cref{eq:right-update}）。这正是 GTSAM 的 \texttt{retract}/\texttt{localCoordinates}、Ceres 的 \texttt{Manifold::Plus/Minus}、Sophus/manif 的右乘更新所封装的接口。

\begin{insight}
一句话主线：\textbf{流形} $\to$ \textbf{切空间局部坐标 $\boldsymbol{\xi}$} $\to$ \textbf{retraction $\mathrm{Exp}$（即 $\boxplus$）} $\to$ 误差线性化后只剩欧氏 $\boldsymbol{\xi}$ $\to$ 套标准 Gauss–Newton/LM $\to$ $\boxplus$ 收缩回流形。本书的\emph{右扰动}约定，就是为了让推导出的雅可比与这些求解器的 retract 约定\emph{逐位对齐}，使手推公式可以直接喂给 GTSAM/Ceres。
\end{insight}

\begin{pitfall}[编程陷阱：误把 $9$（或 $16$）维矩阵当优化维度]
直接把 $\mathbf{R}$ 的 $9$ 个元素交给优化器，会出现“过参数化 + 约束破坏”：维度多了、还得维持正交。$\boxplus$ 的意义就是把局部增量维度压到\emph{真实自由度}（$\SO(3)$ 为 $3$、$\SE(3)$ 为 $6$），从而\emph{无需}约束。在 Ceres/GTSAM 里这对应“给状态注册一个 Manifold/LocalParameterization”，漏掉它是新手最常见的收敛性 bug。
\end{pitfall}

% ════════════════════════════════════════════════════════════════
\section{李群上的不确定度}
\label{sec:uncertainty}

状态估计要带\emph{协方差}。但给 $\mathbf{R}$ 直接加高斯噪声会破坏正交性。Barfoot\cite{barfoot2024state} 的标准做法是：\textbf{把高斯放在切空间（李代数），再用指数映射注入群}。

\paragraph{切空间高斯。}以右扰动定义 $\SE(3)$ 上的高斯：
\begin{equation}
  \mathbf{T}=\bar{\mathbf{T}}\,\mathrm{Exp}(\boldsymbol{\epsilon}),\qquad \boldsymbol{\epsilon}\sim\mathcal{N}(\mathbf{0},\boldsymbol{\Sigma}),\ \boldsymbol{\epsilon}\in\mathbb{R}^6,
  \label{eq:lie-gaussian}
\end{equation}
$\bar{\mathbf{T}}$ 为均值（一个确定位姿），$\boldsymbol{\Sigma}$ 是\emph{切空间里}的协方差。如此 $\mathbf{T}$ 恒在 $\SE(3)$ 上，而“分布”由 $\mathbb{R}^6$ 上的标准高斯诱导。$\SO(3)$ 同理 $\mathbf{R}=\bar{\mathbf{R}}\,\mathrm{Exp}(\boldsymbol{\epsilon})$，$\boldsymbol{\epsilon}\in\mathbb{R}^3$。

\paragraph{位姿复合的协方差传播。}两个相互独立的含噪位姿复合 $\mathbf{T}=\mathbf{T}_1\mathbf{T}_2$，均值为 $\bar{\mathbf{T}}=\bar{\mathbf{T}}_1\bar{\mathbf{T}}_2$，协方差到\emph{一阶}（二阶矩）为
\begin{equation}
  \boldsymbol{\Sigma}\approx\boldsymbol{\Sigma}_1+\mathrm{Ad}(\bar{\mathbf{T}}_1)\,\boldsymbol{\Sigma}_2\,\mathrm{Ad}(\bar{\mathbf{T}}_1)^\top.
  \label{eq:compound-cov}
\end{equation}
这里伴随 $\mathrm{Ad}(\bar{\mathbf{T}}_1)$ 把第二个位姿的切空间协方差“搬运”到复合后的切空间——\textbf{协方差传播的本质，就是用伴随在不同均值点之间搬运切空间协方差}。\cref{eq:compound-cov} 是 VIO 预积分、位姿图边协方差的基石。

\paragraph{$*$ 四阶修正与“香蕉形”分布。}当不确定度较大、或多段位姿链式复合时，一阶传播低估了真实弥散，分布呈现弯曲的“香蕉形”。Barfoot 给出在扰动上展开到\emph{四阶}的修正（见 \cref{sec:lie-appendix} 的指引）：
\begin{equation}
  \boldsymbol{\Sigma}_{4\text{th}}\approx\boldsymbol{\Sigma}_1+\boldsymbol{\Sigma}_2'+\tfrac14\boldsymbol{\mathcal{B}}+\tfrac{1}{12}\big(\boldsymbol{\mathcal{A}}_1\boldsymbol{\Sigma}_2'+\boldsymbol{\Sigma}_2'\boldsymbol{\mathcal{A}}_1^\top+\boldsymbol{\mathcal{A}}_2'\boldsymbol{\Sigma}_1+\boldsymbol{\Sigma}_1\boldsymbol{\mathcal{A}}_2'^\top\big),
  \label{eq:fourth-order}
\end{equation}
其中 $\boldsymbol{\Sigma}_2'=\mathrm{Ad}(\bar{\mathbf{T}}_1)\boldsymbol{\Sigma}_2\mathrm{Ad}(\bar{\mathbf{T}}_1)^\top$，算子 $\boldsymbol{\mathcal{A}},\boldsymbol{\mathcal{B}}$ 的定义见附录指引。本书在状态估计部分（\cref{ch:slam_est}）会再用到；初读可只记一阶式 \cref{eq:compound-cov}。

\begin{insight}
这一节是 Barfoot《状态估计》区别于十四讲的\emph{招牌}：把“位姿带不确定度”讲成“切空间里的高斯 + 伴随搬运”，使协方差传播、位姿求逆、相关位姿差分都有干净闭式。它把第 \ref{sec:adjoint} 节的伴随从“代数技巧”升格为“概率工具”，是滤波/平滑章节的前置。
\end{insight}

% ════════════════════════════════════════════════════════════════
\section{\texorpdfstring{$*$}{*}相似变换群 Sim(3)}
\label{sec:sim3}

单目 SLAM 存在\emph{尺度不确定性与漂移}，$\SE(3)$ 表达不出尺度变化，需显式带一个尺度因子 $s$。相似变换 $\mathbf{p}'=s\mathbf{R}\mathbf{p}+\mathbf{t}$ 构成相似变换群
\begin{equation}
  \mathrm{Sim}(3)=\Big\{\mathbf{S}=\begin{bmatrix}s\mathbf{R}&\mathbf{t}\\\mathbf{0}^\top&1\end{bmatrix}\Big\},\qquad
  \mathfrak{sim}(3)\ni\boldsymbol{\zeta}=[\boldsymbol{\rho}^\top,\boldsymbol{\phi}^\top,\sigma]^\top\in\mathbb{R}^7,
  \label{eq:sim3}
\end{equation}
其李代数多一维尺度 $\sigma$，对应 $s=e^{\sigma}$。指数映射形如 $\exp(\boldsymbol{\zeta}^\wedge)=\big[\begin{smallmatrix}e^{\sigma}\mathbf{R}&\mathbf{J}_s\boldsymbol{\rho}\\\mathbf{0}^\top&1\end{smallmatrix}\big]$，$\mathbf{J}_s$ 的闭式见十四讲\cite{gaoxiang2019slam14}§4.5。$\mathrm{Sim}(3)$ 在单目回环纠正（如 ORB-SLAM 的尺度感知位姿图）中必备；不涉及单目可跳过。

% ════════════════════════════════════════════════════════════════
\section{代码：Sophus 实践}
\label{sec:sophus}

Sophus（基于 Eigen 的头文件库）实现了 $\SO(3),\SE(3),\mathrm{Sim}(3)$。下面演示构造、指数/对数、$\wedge/\vee$ 与\emph{右扰动}更新（注意与本书约定一致地写成\textbf{右乘}）。

\begin{codebox}[C++]{Sophus：构造、exp/log、右扰动更新}
#include <sophus/se3.hpp>
#include <Eigen/Core>
// 沿 Z 轴 90 度
Eigen::Matrix3d R = Eigen::AngleAxisd(M_PI/2, Eigen::Vector3d::UnitZ()).toRotationMatrix();
Sophus::SO3d SO3_R(R);                 // 由旋转矩阵构造（也可 SO3d(q) 由四元数）
Eigen::Vector3d so3 = SO3_R.log();     // 对数映射 -> 李代数(3维)
Eigen::Matrix3d hat = Sophus::SO3d::hat(so3);   // ^ : 向量 -> 反对称阵

// 右扰动更新（本书主线）：X <- X * Exp(delta)
Eigen::Vector3d delta(1e-4, 0, 0);
Sophus::SO3d SO3_updated = SO3_R * Sophus::SO3d::exp(delta);   // 右乘！

// SE(3) 同理；注意 Sophus 的 se(3) 顺序是“平移在前、旋转在后”，与本书一致
Eigen::Vector3d t(1, 0, 0);
Sophus::SE3d SE3(R, t);
Eigen::Matrix<double,6,1> xi = SE3.log();        // [rho; phi]
Eigen::Matrix<double,6,1> dxi; dxi.setZero(); dxi(0)=1e-4;
Sophus::SE3d SE3_updated = SE3 * Sophus::SE3d::exp(dxi);       // 右乘更新
\end{codebox}

\begin{note}
十四讲原书代码用\emph{左}乘 $\mathrm{Exp}(\delta)\,\mathbf{X}$；本书统一改为\emph{右}乘 $\mathbf{X}\,\mathrm{Exp}(\delta)$ 以对齐右扰动约定。容器存放需对齐分配器：\texttt{std::vector<Sophus::SE3d, Eigen::aligned\_allocator<Sophus::SE3d>>}。轨迹误差的核心是一行 \texttt{(T\_gt.inverse()*T\_est).log().norm()}，即 $\|\mathrm{Log}(\bar{\mathbf{T}}_{\text{gt}}^{-1}\mathbf{T}_{\text{est}})\|$——\emph{用李代数把位姿误差变成可求范数的向量}。
\end{note}

% ════════════════════════════════════════════════════════════════
\section{本章推导附录}
\label{sec:lie-appendix}

本章主要证明（$\so(3)$ 的引出、Rodrigues 公式、左雅可比 $\mathbf{J}$、左/右扰动求导）均已在正文相应定理处给出。此处仅收录两处\emph{极冗长}的补充闭式，供正文按需取用。

\begin{derivation}[$\SE(3)$ 的 $\mathbf{Q}_l$ 块闭式（Barfoot 式 8.91）]
\cref{eq:se3-6x6} 中 $6\times6$ 左雅可比的耦合块 $\mathbf{Q}_l$ 由双重级数 $\mathbf{Q}_l=\sum_{n,m\ge0}\frac{(\boldsymbol{\phi}^\wedge)^n\boldsymbol{\rho}^\wedge(\boldsymbol{\phi}^\wedge)^m}{(n+m+2)!}$ 归并而得（$\theta=\|\boldsymbol{\phi}\|$）：
\begin{align*}
  \mathbf{Q}_l(\boldsymbol{\rho},\boldsymbol{\phi})
  &=\tfrac12\boldsymbol{\rho}^\wedge
   +\frac{\theta-\sin\theta}{\theta^3}\big(\boldsymbol{\phi}^\wedge\boldsymbol{\rho}^\wedge+\boldsymbol{\rho}^\wedge\boldsymbol{\phi}^\wedge+\boldsymbol{\phi}^\wedge\boldsymbol{\rho}^\wedge\boldsymbol{\phi}^\wedge\big)\\
  &\quad+\frac{\theta^2+2\cos\theta-2}{2\theta^4}\big(\boldsymbol{\phi}^\wedge\boldsymbol{\phi}^\wedge\boldsymbol{\rho}^\wedge+\boldsymbol{\rho}^\wedge\boldsymbol{\phi}^\wedge\boldsymbol{\phi}^\wedge-3\,\boldsymbol{\phi}^\wedge\boldsymbol{\rho}^\wedge\boldsymbol{\phi}^\wedge\big)\\
  &\quad+\frac{2\theta-3\sin\theta+\theta\cos\theta}{2\theta^5}\big(\boldsymbol{\phi}^\wedge\boldsymbol{\rho}^\wedge\boldsymbol{\phi}^\wedge\boldsymbol{\phi}^\wedge+\boldsymbol{\phi}^\wedge\boldsymbol{\phi}^\wedge\boldsymbol{\rho}^\wedge\boldsymbol{\phi}^\wedge\big).
\end{align*}
右版本 $\mathbf{Q}_r(\boldsymbol{\xi})=\mathbf{Q}_l(-\boldsymbol{\xi})$。常规 SLAM 优化旋转较小，取 $\boldsymbol{\mathcal{J}}\approx\mathbf{I}$ 即可；仅高精度协方差传播时用此式。
\end{derivation}

\begin{derivation}[位姿复合四阶协方差的算子（\cref{eq:fourth-order}）]
\cref{eq:fourth-order} 中的算子定义为（Barfoot\cite{barfoot2024state} §8.3，$(\cdot)^\curlywedge$ 为 \cref{eq:ad} 的小伴随）：
\[
  \boldsymbol{\mathcal{A}}_i=\mathbb{E}\!\big[\boldsymbol{\epsilon}_i^\curlywedge\boldsymbol{\epsilon}_i^\curlywedge\big],\qquad
  \boldsymbol{\mathcal{B}}=\mathbb{E}\!\big[\boldsymbol{\epsilon}_1^\curlywedge\,(\boldsymbol{\epsilon}_2'\boldsymbol{\epsilon}_2'^\top)\,\boldsymbol{\epsilon}_1^{\curlywedge\top}\big],\qquad
  \boldsymbol{\epsilon}_2'=\mathrm{Ad}(\bar{\mathbf{T}}_1)\boldsymbol{\epsilon}_2.
\]
该修正在扰动上为四阶、在协方差上为二阶；当两位姿不相关时与 sigma-point（无迹）法代数等价。
\end{derivation}

% ════════════════════════════════════════════════════════════════
\section*{本章常见误解汇总}
\begin{table}[H]\centering\small
\begin{tabular}{p{0.46\textwidth} p{0.46\textwidth}}
\toprule
\textbf{常见误解} & \textbf{正确理解} \\
\midrule
$\SO(3)$ 是向量空间，可以加减 & 只对乘法封闭，是群非向量空间（\cref{eq:not-vs}）\\
$\mathbf{R}^{-1}=\mathbf{R}^\top$ 是巧合 & 正交性的定义性质，旋转保范的必然 \\
$\se(3)$ 的 $\boldsymbol{\rho}$ 就是平移 $\mathbf{t}$ & $\mathbf{t}=\mathbf{J}\boldsymbol{\rho}$，仅 $\boldsymbol{\phi}\to0$ 时相等（\cref{thm:se3-exp}）\\
指数映射一一对应 & 满射非单射，需限 $|\theta|<\pi$（\cref{prop:surjective}）\\
BCH 是精确等式 & 仅小扰动下的一阶近似（\cref{eq:bch-linear}）\\
左右扰动可随意混用 & 雅可比/更新/协方差须同侧，否则发散 \\
$\mathrm{Ad}(\mathbf{T})$ 是 $4\times4$ 位姿 & 是 $6\times6$、作用于切向量 $\boldsymbol{\xi}$ 的线性映射 \\
直接优化 $\mathbf{R}$ 的 $9$ 元素 & 过参数化 + 破坏约束；应用 $\boxplus$ 在 $3/6$ 维切空间更新 \\
\bottomrule
\end{tabular}
\end{table}

% ════════════════════════════════════════════════════════════════
\section*{本章小结}

\begin{table}[H]\centering\small
\caption{符号表}
\begin{tabular}{lll}
\toprule
符号 & 含义 & 首次出现 \\
\midrule
$\SO(3),\SE(3)$ & 三维旋转群 / 位姿群 & \cref{eq:so3,eq:se3} \\
$\so(3),\se(3)$ & 对应李代数 & \cref{def:so3-se3} \\
$(\cdot)^\wedge,(\cdot)^\vee$ & 反对称化 / 其逆（$\se(3)$ 为重载） & \cref{eq:hat} \\
$\mathrm{Exp},\mathrm{Log}$ & 吃/吐向量的指数 / 对数 & \cref{sec:exp-log} \\
$\boldsymbol{\phi}=\theta\mathbf{a},\ \boldsymbol{\xi}=[\boldsymbol{\rho};\boldsymbol{\phi}]$ & 轴角 / $\se(3)$ 坐标（平移在前） & \cref{def:so3-se3} \\
$\mathbf{J}_l,\mathbf{J}_r$ & 左 / 右雅可比（$3\times3$） & \cref{eq:Jl,eq:Jr} \\
$\boldsymbol{\mathcal{J}}_l,\mathbf{Q}_l$ & $\SE(3)$ 的 $6\times6$ 雅可比及其耦合块 & \cref{eq:se3-6x6} \\
$\mathrm{Ad}(\mathbf{T}),\ \boldsymbol{\xi}^\curlywedge$ & 群伴随（$6\times6$）/ 小伴随 & \cref{eq:Ad,eq:ad} \\
$(\cdot)^\odot$ & 齐次点 $\to 4\times6$ 的扰动算子 & \cref{eq:right-pert-se3} \\
$\boxplus,\boxminus$ & 流形“加/减”（retraction/局部坐标） & \cref{eq:boxplus} \\
$\bar{\mathbf{T}},\boldsymbol{\Sigma}$ & 位姿均值 / 切空间协方差 & \cref{eq:lie-gaussian} \\
\bottomrule
\end{tabular}
\end{table}

\begin{table}[H]\centering\small
\caption{定理速查表}
\begin{tabular}{lll}
\toprule
定理 / 公式 & 一句话说明 & 对应 \\
\midrule
罗德里格斯 \cref{eq:rodrigues} & 轴角 $\to$ 旋转矩阵，避免无穷级数 & \cref{thm:so3-exp} \\
$\SE(3)$ 指数 \cref{eq:left-jacobian} & $\mathbf{t}=\mathbf{J}\boldsymbol{\rho}$，左雅可比 $\mathbf{J}$ 现身 & \cref{thm:se3-exp} \\
BCH 线性近似 \cref{eq:bch-linear} & 李代数加法 $\leftrightarrow$ 李群乘法，分左右 & \cref{sec:bch} \\
右扰动求导 \cref{eq:right-pert-so3} & $\partial(\mathbf{R}\mathbf{p})/\partial\boldsymbol{\varphi}=-\mathbf{R}\mathbf{p}^\wedge$，无 $\mathbf{J}$ & \cref{sec:perturb} \\
伴随移指数 \cref{eq:Ad-move} & 左右扰动精确互换的桥 & \cref{sec:adjoint} \\
$\boxplus/\boxminus$ \cref{eq:boxplus} & 流形优化 = 切空间欧氏优化 + 收缩 & \cref{sec:manifold} \\
复合协方差 \cref{eq:compound-cov} & 伴随搬运切空间协方差 & \cref{sec:uncertainty} \\
\bottomrule
\end{tabular}
\end{table}

\begin{table}[H]\centering\small
\caption{知识点总表}
\begin{tabular}{clll}
\toprule
\# & 知识点 & 核心要点 & 难度 \\
\midrule
1 & 群与矩阵李群 & 加法不封闭、乘法封闭；封结幺逆 & \(\star\star\) \\
2 & 李代数 & $\mathbf{R}\mathbf{R}^\top=\mathbf{I}$ 求导挤出反对称 & \(\star\star\star\) \\
3 & 指数/对数 & Rodrigues、$\SE(3)$ 指数、$\mathbf{J}$ & \(\star\star\star\) \\
4 & BCH 与 $\mathbf{J}_l,\mathbf{J}_r$ & 李代数加 $\leftrightarrow$ 群乘的线性近似 & \(\star\star\star\star\) \\
5 & 右扰动模型 & $-\mathbf{R}\mathbf{p}^\wedge$；更新右乘指数 & \(\star\star\star\star\) \\
6 & 伴随 $\mathrm{Ad},\mathrm{ad}$ & 左右扰动/坐标系/协方差的桥 & \(\star\star\star\star\) \\
7 & 流形/$\boxplus\boxminus$ & 对接因子图、GTSAM/Ceres & \(\star\star\star\star\) \\
8 & 李群不确定度 & 切空间高斯 + 伴随传播 & \(\star\star\star\star\star\) \\
\bottomrule
\end{tabular}
\end{table}

% ════════════════════════════════════════════════════════════════
\section*{延伸阅读}
\begin{itemize}
  \item \textbf{严谨主线}：Barfoot《State Estimation for Robotics》\cite{barfoot2024state} 第 7--8 章——矩阵李群、$6\times6$ 雅可比、李群上的不确定度与位姿复合（本章不确定度一节的完整出处）。
  \item \textbf{几何直觉与代码}：视觉 SLAM 十四讲\cite{gaoxiang2019slam14} 第 4 章——招牌的“从 $\mathbf{R}\mathbf{R}^\top=\mathbf{I}$ 引出李代数”、Sophus 实践。
  \item \textbf{现代统一记号}：Solà 等《A micro Lie theory for state estimation in robotics》\cite{sola2018micro}——右/左 $\oplus\,\ominus$、函数的右雅可比与链式法则，强烈推荐对照本书右扰动约定阅读。
  \item \textbf{流形优化接口}：Hertzberg 等\cite{hertzberg2013integrating}（$\boxplus/\boxminus$ 封装的原始论文）、SLAM Handbook\cite{carlone2026handbook} 第 2 章（流形/retraction/因子图衔接）。
  \item \textbf{随机李群理论}：Chirikjian\cite{chirikjian2011stochastic}——李群上概率与信息论的奠基性参考。
\end{itemize}

\section*{本章与后续章节的关系}
\begin{table}[H]\centering\small
\begin{tabular}{lll}
\toprule
后续章节 & 关系 & 本章铺垫 \\
\midrule
\cref{ch:nlopt} 非线性优化 & 位姿作为流形变量做 GN/LM & 右扰动雅可比、$\boxplus$ 更新 \\
\cref{ch:slam_est} 状态估计 & 携带协方差的位姿估计 & 切空间高斯、伴随传播 \\
\cref{ch:vo} 视觉里程计 & 重投影误差对位姿求导 & \cref{eq:right-pert-se3} 的 $\mathbf{T}\tilde{\mathbf{p}}^\odot$ \\
\cref{ch:vio} VIO & IMU 预积分在 $\SE(3)$/$\SO(3)$ 上 & 伴随、$\SE(3)$ 雅可比、不确定度 \\
\bottomrule
\end{tabular}
\end{table}

\section*{故障排查手册}
\begin{enumerate}
  \item \textbf{症状}：BA / 位姿图\emph{不收敛或发散}。\textbf{排查}：检查雅可比扰动侧与更新乘法侧是否同为右侧（\cref{sec:perturb} 陷阱）；检查是否给求解器注册了 $\SO(3)/\SE(3)$ 的流形/局部参数化（\cref{sec:manifold} 陷阱）。
  \item \textbf{症状}：大旋转处残差/雅可比\emph{突变}。\textbf{原因}：$\mathrm{Log}$ 在 $|\theta|\to\pi$ 的 $2\pi$ 跳变（\cref{prop:surjective}）。\textbf{对策}：对相对旋转取 $\mathrm{Log}$、保持每步增量小。
  \item \textbf{症状}：协方差传播后\emph{不对称或非正定}。\textbf{排查}：确认用 $\mathrm{Ad}(\bar{\mathbf{T}}_1)\boldsymbol{\Sigma}_2\mathrm{Ad}(\bar{\mathbf{T}}_1)^\top$（\cref{eq:compound-cov}）而非漏掉转置；大不确定度需四阶修正（\cref{eq:fourth-order}）。
  \item \textbf{症状}：$\mathbf{J}$ 或 $\mathbf{J}^{-1}$ 计算出 \texttt{NaN}。\textbf{原因}：$\theta\to0$ 时 $\sin\theta/\theta$、$\theta\cot(\theta/2)$ 的 $0/0$。\textbf{对策}：小角度用泰勒展开分支（$\mathbf{J}\approx\mathbf{I}+\tfrac12\boldsymbol{\phi}^\wedge$）。
  \item \textbf{症状}：Sophus 里 $\se(3)$ 分量顺序与公式对不上。\textbf{原因}：Sophus/本书为“平移在前、旋转在后”，部分文献（JPL、某些 GTSAM 约定）相反。\textbf{对策}：核对 $\boldsymbol{\xi}=[\boldsymbol{\rho};\boldsymbol{\phi}]$ 顺序后再套雅可比。
\end{enumerate}
